\documentclass[12pt,a4paper]{article}

\usepackage[brazil]{babel}
\usepackage[utf8]{inputenc}
\usepackage[T1]{fontenc}
\usepackage[a4paper,margin=2.0cm]{geometry}
\usepackage{amsmath,amssymb}
\usepackage{enumitem}
\usepackage{circuitikz}
\usepackage{xcolor}

\begin{document}

\begin{center}
{\LARGE \textbf{Lista de Exercícios}}\\[0.2cm]
{\Large Resistência Elétrica, Associação de Resistores e 2ª Lei de Ohm}
\end{center}

\vspace{0.5cm}

\section*{Introdução Teórica}

A resistência elétrica é uma propriedade dos materiais que mede a oposição à passagem da corrente elétrica. Sua unidade no Sistema Internacional é o ohm ($\Omega$).

A relação fundamental entre tensão, corrente e resistência é dada pela Primeira Lei de Ohm:

\[
V = RI
\]

onde:

\begin{itemize}
\item $V$ é a tensão elétrica (V);
\item $I$ é a corrente elétrica (A);
\item $R$ é a resistência elétrica ($\Omega$).
\end{itemize}

\subsection*{Associação em Série}

Quando os resistores são ligados em série:

\begin{itemize}
\item a corrente é a mesma em todos os elementos;
\item a resistência equivalente é dada por
\end{itemize}

\[
R_{eq}=R_1+R_2+R_3+\cdots
\]

\subsection*{Associação em Paralelo}

Quando os resistores são ligados em paralelo:

\begin{itemize}
\item a tensão é a mesma em todos os resistores;
\item a resistência equivalente é calculada por
\end{itemize}

\[
\frac{1}{R_{eq}}
=
\frac{1}{R_1}
+
\frac{1}{R_2}
+
\frac{1}{R_3}
+\cdots
\]

\subsection*{Associação Mista}

Circuitos que apresentam simultaneamente ligações em série e em paralelo são chamados de circuitos mistos. A determinação da resistência equivalente é feita simplificando o circuito por etapas.

\subsection*{Segunda Lei de Ohm}

A resistência de um fio homogêneo pode ser calculada por

\[
R=\rho\frac{L}{A}
\]

onde:

\begin{itemize}
\item $\rho$ é a resistividade do material;
\item $L$ é o comprimento do fio;
\item $A$ é a área da seção transversal.
\end{itemize}

\newpage

\section*{Exercícios}

\begin{enumerate}[label=\textbf{\arabic*.}]

%%%%%%%%%%%%%%%%%%%%%%%%%%%%%%%%%%%%%%%%%%%%%%%%%%%%%%%%%%%%

\item Determine a resistência equivalente do circuito.

\begin{center}
\begin{circuitikz}
\draw
(0,0) to[R=$10\Omega$] (2,0)
      to[R=$20\Omega$] (4,0);
\end{circuitikz}
\end{center}

\vspace{1cm}

%%%%%%%%%%%%%%%%%%%%%%%%%%%%%%%%%%%%%%%%%%%%%%%%%%%%%%%%%%%%

\item Determine a resistência equivalente do circuito.

\begin{center}
\begin{circuitikz}
\draw
(0,0) to[R=$15\Omega$] (2,0)
      to[R=$25\Omega$] (4,0)
      to[R=$10\Omega$] (6,0);
\end{circuitikz}
\end{center}

\vspace{1cm}

%%%%%%%%%%%%%%%%%%%%%%%%%%%%%%%%%%%%%%%%%%%%%%%%%%%%%%%%%%%%

\item Uma fonte de $24\,V$ alimenta o circuito abaixo. Determine a corrente total.

\begin{center}
\begin{circuitikz}
\draw
(0,0) to[battery1,l=$24V$] (0,3)
      -- (1,3)
      to[R=$8\Omega$] (3,3)
      to[R=$4\Omega$] (5,3)
      -- (5,0)
      -- (0,0);
\end{circuitikz}
\end{center}

\vspace{1cm}

%%%%%%%%%%%%%%%%%%%%%%%%%%%%%%%%%%%%%%%%%%%%%%%%%%%%%%%%%%%%

\item Determine a resistência equivalente.

\begin{center}
\begin{circuitikz}
\draw (0,2)--(4,2);
\draw (0,0)--(4,0);

\draw (1,2) to[R=$12\Omega$] (1,0);
\draw (3,2) to[R=$6\Omega$] (3,0);
\end{circuitikz}
\end{center}

\vspace{1cm}

%%%%%%%%%%%%%%%%%%%%%%%%%%%%%%%%%%%%%%%%%%%%%%%%%%%%%%%%%%%%

\item Determine a resistência equivalente.

\begin{center}
\begin{circuitikz}
\draw (0,2)--(6,2);
\draw (0,0)--(6,0);

\draw (1,2) to[R=$10\Omega$] (1,0);
\draw (3,2) to[R=$15\Omega$] (3,0);
\draw (5,2) to[R=$30\Omega$] (5,0);
\end{circuitikz}
\end{center}

\vspace{1cm}

%%%%%%%%%%%%%%%%%%%%%%%%%%%%%%%%%%%%%%%%%%%%%%%%%%%%%%%%%%%%

\item Uma fonte de $12V$ alimenta dois resistores de $6\Omega$ em paralelo.

Determine:

\begin{enumerate}
\item resistência equivalente;
\item corrente total.
\end{enumerate}

\begin{center}
\begin{circuitikz}
\draw
(0,0) to[battery1,l=$12V$] (0,4)
      -- (4,4)
      -- (4,0)
      -- (0,0);

\draw (1,4) to[R=$6\Omega$] (1,0);
\draw (3,4) to[R=$6\Omega$] (3,0);
\end{circuitikz}
\end{center}

\vspace{1cm}

%%%%%%%%%%%%%%%%%%%%%%%%%%%%%%%%%%%%%%%%%%%%%%%%%%%%%%%%%%%%

\item Determine a resistência equivalente.

\begin{center}
\begin{circuitikz}
\draw
(0,0) to[R=$5\Omega$] (2,0)
      -- (4,0);

\draw (2,0)--(2,2);
\draw (4,0)--(4,2);

\draw (2,2) to[R=$10\Omega$] (4,2);
\draw (2,0) to[R=$10\Omega$] (4,0);
\end{circuitikz}
\end{center}

\vspace{1cm}

%%%%%%%%%%%%%%%%%%%%%%%%%%%%%%%%%%%%%%%%%%%%%%%%%%%%%%%%%%%%

\item Determine a resistência equivalente.

\begin{center}
\begin{circuitikz}
\draw
(0,0) to[R=$8\Omega$] (2,0);

\draw (2,0)--(2,2);
\draw (4,0)--(4,2);

\draw (2,2) to[R=$12\Omega$] (4,2);
\draw (2,0) to[R=$6\Omega$] (4,0);
\end{circuitikz}
\end{center}

\vspace{1cm}

%%%%%%%%%%%%%%%%%%%%%%%%%%%%%%%%%%%%%%%%%%%%%%%%%%%%%%%%%%%%

\item Calcule a corrente fornecida por uma fonte de $30V$.

\begin{center}
\begin{circuitikz}
\draw
(0,0) to[battery1,l=$30V$] (0,3)
      -- (2,3)
      to[R=$10\Omega$] (4,3)
      to[R=$5\Omega$] (6,3)
      -- (6,0)
      -- (0,0);
\end{circuitikz}
\end{center}

\vspace{1cm}

%%%%%%%%%%%%%%%%%%%%%%%%%%%%%%%%%%%%%%%%%%%%%%%%%%%%%%%%%%%%

\item Determine a resistência equivalente.

\begin{center}
\begin{circuitikz}
\draw
(0,0) to[R=$4\Omega$] (2,0);

\draw (2,0)--(2,2);
\draw (5,0)--(5,2);

\draw (2,2) to[R=$6\Omega$] (5,2);
\draw (2,0) to[R=$3\Omega$] (5,0);

\draw (5,0) to[R=$2\Omega$] (7,0);
\end{circuitikz}
\end{center}

\vspace{1cm}

%%%%%%%%%%%%%%%%%%%%%%%%%%%%%%%%%%%%%%%%%%%%%%%%%%%%%%%%%%%%

\item Determine a resistência equivalente.

\begin{center}
\begin{circuitikz}
\draw
(0,0) to[R=$2\Omega$] (2,0)
      to[R=$3\Omega$] (4,0)
      to[R=$5\Omega$] (6,0);
\end{circuitikz}
\end{center}

\vspace{1cm}

%%%%%%%%%%%%%%%%%%%%%%%%%%%%%%%%%%%%%%%%%%%%%%%%%%%%%%%%%%%%

\item Determine a resistência equivalente.

\begin{center}
\begin{circuitikz}
\draw (0,2)--(5,2);
\draw (0,0)--(5,0);

\draw (1,2) to[R=$4\Omega$] (1,0);
\draw (3,2) to[R=$8\Omega$] (3,0);
\draw (5,2) to[R=$8\Omega$] (5,0);
\end{circuitikz}
\end{center}

\vspace{1cm}

%%%%%%%%%%%%%%%%%%%%%%%%%%%%%%%%%%%%%%%%%%%%%%%%%%%%%%%%%%%%

\item Determine a resistência equivalente.

\begin{center}
\begin{circuitikz}
\draw
(0,0) to[R=$12\Omega$] (2,0);

\draw (2,0)--(2,2);
\draw (5,0)--(5,2);

\draw (2,2) to[R=$18\Omega$] (5,2);
\draw (2,0) to[R=$9\Omega$] (5,0);

\draw (5,0) to[R=$6\Omega$] (7,0);
\end{circuitikz}
\end{center}

\vspace{1cm}

%%%%%%%%%%%%%%%%%%%%%%%%%%%%%%%%%%%%%%%%%%%%%%%%%%%%%%%%%%%%

\item Um fio de cobre possui:

\[
L=20\,m
\]

\[
A=1,5\times10^{-6}\,m^2
\]

\[
\rho=1,7\times10^{-8}\,\Omega\cdot m
\]

Determine sua resistência.

\vspace{2cm}

%%%%%%%%%%%%%%%%%%%%%%%%%%%%%%%%%%%%%%%%%%%%%%%%%%%%%%%%%%%%

\item Um fio de alumínio possui:

\[
L=50\,m
\]

\[
A=2,0\times10^{-6}\,m^2
\]

\[
\rho=2,8\times10^{-8}\,\Omega\cdot m
\]

Determine sua resistência.

\vspace{2cm}

%%%%%%%%%%%%%%%%%%%%%%%%%%%%%%%%%%%%%%%%%%%%%%%%%%%%%%%%%%%%

\item Um fio possui resistência de $8\Omega$.

Se o comprimento for duplicado e a área permanecer constante, determine a nova resistência.

\vspace{2cm}

%%%%%%%%%%%%%%%%%%%%%%%%%%%%%%%%%%%%%%%%%%%%%%%%%%%%%%%%%%%%

\item Um fio possui resistência de $12\Omega$.

Sua área transversal é duplicada sem alterar o comprimento.

Determine a nova resistência.

\vspace{2cm}

%%%%%%%%%%%%%%%%%%%%%%%%%%%%%%%%%%%%%%%%%%%%%%%%%%%%%%%%%%%%

\item Determine a resistência equivalente.

\begin{center}
\begin{circuitikz}
\draw
(0,0) to[R=$5\Omega$] (2,0);

\draw (2,0)--(2,2);
\draw (5,0)--(5,2);

\draw (2,2) to[R=$20\Omega$] (5,2);
\draw (2,0) to[R=$20\Omega$] (5,0);

\draw (5,0) to[R=$10\Omega$] (7,0);
\end{circuitikz}
\end{center}

\vspace{1cm}

%%%%%%%%%%%%%%%%%%%%%%%%%%%%%%%%%%%%%%%%%%%%%%%%%%%%%%%%%%%%

\item Uma fonte de $48V$ alimenta o circuito abaixo.

Determine a corrente total.

\begin{center}
\begin{circuitikz}
\draw
(0,0) to[battery1,l=$48V$] (0,4)
      -- (2,4);

\draw
(2,4) to[R=$12\Omega$] (2,0);

\draw
(4,4) to[R=$12\Omega$] (4,0);

\draw (2,4)--(4,4);
\draw (2,0)--(4,0);

\draw (4,0)--(0,0);
\end{circuitikz}
\end{center}

\vspace{1cm}

%%%%%%%%%%%%%%%%%%%%%%%%%%%%%%%%%%%%%%%%%%%%%%%%%%%%%%%%%%%%

\item Determine a resistência equivalente entre os terminais A e B.

\begin{center}
\begin{circuitikz}
\draw
(0,0) node[left]{A}
      to[R=$4\Omega$] (2,0);

\draw (2,0)--(2,2);
\draw (5,0)--(5,2);

\draw (2,2) to[R=$8\Omega$] (5,2);
\draw (2,0) to[R=$8\Omega$] (5,0);

\draw
(5,0) to[R=$4\Omega$] (7,0)
      node[right]{B};
\end{circuitikz}
\end{center}

\vspace{1cm}

\end{enumerate}

\newpage

\section*{Desafio}

Determine a resistência equivalente do circuito abaixo sem utilizar calculadora.

\begin{center}
\begin{circuitikz}
\draw
(0,0) to[R=$6\Omega$] (2,0);

\draw (2,0)--(2,2);
\draw (6,0)--(6,2);

\draw
(2,2) to[R=$12\Omega$] (4,2)
      to[R=$12\Omega$] (6,2);

\draw
(2,0) to[R=$24\Omega$] (6,0);

\draw
(6,0) to[R=$6\Omega$] (8,0);
\end{circuitikz}
\end{center}

\vfill

\begin{center}
\textbf{Boa resolução!}
\end{center}

\end{document}